\documentclass[a4paper,11pt]{letter}
\usepackage[utf8]{inputenc}
\usepackage[T1]{fontenc}
\usepackage[french]{babel}
\usepackage{geometry}
\geometry{left=2.5cm,right=2.5cm,top=2cm,bottom=2cm}
\usepackage{hyperref}
\hypersetup{colorlinks=true, urlcolor=blue}

\begin{document}

\begin{flushleft}
    \textbf{Ismail AISSAOUI} \\
    Ingénieur d'État en Intelligence Artificielle (ESI-SBA) \\
    Email: ismail.aissaoui.pro@gmail.com \\
    Téléphone: +213 660 70 77 96 \\
    LinkedIn: https://linkedin.com/in/aissaoui-ismail-6a77a92a7 \\
    Portfolio: https://ismail-aissaoui.vercel.app
\end{flushleft}

\bigskip

\begin{flushright}
    À l'attention de Monsieur Johann Bourcier \\
    UPPA / Inria \\
    Anglet, France
\end{flushright}

\bigskip

\begin{center}
    \textbf{Objet : Candidature au doctorat : Composition et configuration d'architectures de jumeaux numériques (EDT)}
\end{center}

\bigskip

Monsieur le Docteur Bourcier,

Ingénieur d'État en Intelligence Artificielle diplômé de l'ESI-SBA, je souhaite vous soumettre ma candidature pour le poste de doctorat intitulé « Composing and configuring digital twin architectures to meet user requirements » au sein du programme national Engineering Digital Twin (EDT). Cette candidature est motivée par la forte adéquation de mon profil avec vos recherches sur l'optimisation des architectures logicielles et l'exploration d'espaces de conception (DSE).

Le défi de naviguer dans un espace de conception architectural pour équilibrer des exigences non-fonctionnelles concurrentes (énergie, performance, budget) est au cœur de mes travaux actuels. Dans ma thèse de fin d'études chez Sonatrach, j'ai conçu un **Jumeau Numérique informé par la physique** pour la surveillance de turbines à gaz où j'ai été confronté au « paradoxe de la latence ». Pour le résoudre, j'ai mis en œuvre une stratégie de co-conception matériel-logiciel, optimisant des architectures hybrides **Mamba-SSM** et **xLSTM** pour atteindre une inférence de **0,04 ms** sur du matériel contraint. Cette expérience m'a permis de comprendre comment les choix architecturaux impactent directement les contraintes opérationnelles.

Je possède une solide maîtrise des algorithmes métaheuristiques (PSO, Algorithmes Génétiques) et des techniques d'optimisation bayésienne, outils que je considère essentiels pour automatiser la configuration des architectures de jumeaux numériques. Ma capacité à manipuler des outils comme PyTorch et C++ JIT, combinée à mon intérêt pour les systèmes distribués, me permettra de contribuer efficacement au développement de prototypes au sein du laboratoire LIUPPA.

Rejoindre votre équipe au sein du projet EDT représente pour moi une opportunité unique de transformer mes compétences en optimisation de systèmes en contributions académiques de premier plan. Je suis convaincu que mon approche pragmatique et mon expertise technique sauront soutenir les objectifs de votre groupe de recherche.

Je vous remercie de l'attention que vous porterez à ma candidature et reste à votre disposition pour un entretien.

Je vous prie d'agréer, Monsieur le Docteur, l'expression de mes salutations distinguées.

\bigskip

Ismail AISSAOUI

\end{document}
